\documentclass[11pt]{article}

\usepackage{jcd-macros}
\usepackage[numbers,square]{natbib}
\bluehyperref  % Dark blue hyperref
\usemedgeometry  % Use medium geometry (3cm left/right margin, 2.54 top/bottom)

\begin{document}

\title{26-12 Multi-layer Network Clustering}

\maketitle

\section{Some minor notes on discussion}

\paragraph{Responsible Parties:} Lucas

I didn't take any notes: Lucas, can you write a paragraph?

\section{Lucas' contribution}
My suggestion was to consider a single graph with multidimensional edges, i.e., each edge is a vector of dimension $\ell$ and each entry corresponds to that edge in a given layer. Equivalently we could think of this as an adjacency 3D tensor. Then the question is whether there's a principled way to collapse the tensor into the usual 2D adjacency matrix, at which point standard clustering algorithms can be applied. 

There are many simple ways to do this collapsing (e.g., each edge vector could be collapsed to its norm), but Tatsu suggested learning a good collapsing function from the labeled nodes, which sounded promising but maybe he can elaborate on?

\end{document}
